
\section{Related Work}

The automated fact-checking process entails a three-stage pipeline namely claim detection, evidence retrieval and claim-verification. In this section, we discuss limitations of prior datasets and motivate the need for a realistic open-domain setup currently absent in existing datasets.

Majority of the existing datasets primarily focus on synthetically curated claims from Wikipedia articles \cite{jiang2020hover,thorne-etal-2018-fever,aly2021feverous,creditassess,schuster2021vitamin,sathe-etal-2020-automated}. Evaluation of automated fact-checking systems on such benchmarks is not representative of real-world claims encountered by fact-checkers motivating the need for curation of such real-world claims.  Recent efforts like \cite{kamoi2023wice} focus on creating more realistic claims from Wikipedia by identifying cited statements, but fall short and do not reflect the typical distribution of claims verified by fact-checkers as the claims are still synthetic. More recently, substantial effort has been diverted to collection of previously fact-checked real-world claims in politics \cite{multifc,averitec,claimdecomp}, science \cite{wadden-etal-2020-fact,vladika2023scientific,wright2022generating}, health \cite{kotonya2020explainable} and climate \cite{diggelmann2021climatefever} domains.  However, none of these benchmarks primarily focus on numerical claims, which comprise a substantial portion of real-world claims occurring in political debates or other contexts which comprise numerical quantities. Hence, automated fact-checking systems must account for numerical understanding capabilities that are required for evidence retrieval and claim verification to fact-check such claims. Moreover, existing datasets also do not have a realistic evidence collection due to leakage of manual fact-checker ``gold" evidence documents and related verdicts. This leakage may induce the claim verification models to learn shortcuts or overfit to fact-verification justification from manual fact-checkers. Additionally, the data collection pipeline for these benchmarks rely on heuristic search from open web search engines which solely uses the claim as search query. While this may work well for simple synthetic claims, real-world claims are often complex that require multiple aspects to be verified. Hence, using only the original claims as query may limit the evidences collected due to reliance on surface level matches to claim ignoring the numerical and other aspects of the claim.


More recently, QuanTemp \cite{quantemp} addresses these gaps by curating a realistic open-domain benchmark focusing entirely on numerical claims. While QuanTemp addresses leakage of ``gold" evidence documents by filtering out fact-checking websites, it still suffers from \textit{temporal leakage} issue. This issue occurs when evidence snippets published after the date the claim to be verified was made are also considered as relevant evidence and included in the collection. While Averitec \cite{averitec} proposes to mitigate temporal leakage, it neither focuses on numerical claims nor on handling leakage of ``gold" evidence documents. To mitigate these issues, we handle both temporal and ``gold" evidence leakage. Additionally, to cover all aspects of claims, we propose a claim decomposition process which emulates the fact-checker search process. We achieve this by using their justification description as part of data collection pipeline with particular focus on generating queries around numerical aspects of the claims. This results in high quality relevant evidence compared to existing approaches.